\section{The velocity loop and the drift observer}
\label{sec:vlaw}

Nav2 sends a body velocity; the boat's velocity law turns it into effort
(Figure~\ref{fig:loop}). Whatever this loop does between the command and the velocity is what
MPPI does not model, so it is tuned first and read by every later rule.

\begin{figure}[tbp]
\centering
\begin{tikzpicture}[>=Stealth, font=\footnotesize,
  blk/.style={draw, thick, rounded corners=2pt, fill=white, align=center, inner sep=3pt,
    minimum height=1.25cm},
  sub/.style={font=\scriptsize, sD!40!black},
  sig/.style={->, thick},
  dist/.style={->, thick, sB, decorate, decoration={snake, amplitude=1pt, segment length=5pt,
    post length=4pt}},
  note/.style={font=\scriptsize, align=center}]
  % ---- top row: command path ---------------------------------------------------------
  \node[blk, fill=sA!10, text width=1.9cm] (mppi) at (0,0)
    {Nav2 MPPI\\[1pt] {\scriptsize plans as if\\ $v = $ command}};
  \node[blk, text width=1.5cm, right=0.85cm of mppi] (smooth) {velocity\\smoother};
  \node[blk, text width=2.0cm, right=0.45cm of smooth] (mux)
    {\code{command\_mux}\\[1pt] {\scriptsize hold $\Rightarrow$ neutral}};
  \node[blk, text width=6.3cm, right=0.45cm of mux, align=left] (law)
    {\textbf{velocity law}\\[1pt]
     feed-forward through the speed curves $e|e| = d_1 v + d_2 v|v|$\\
     \quad at $v = v^* - b - R^\top\dhat$ \ (fig.~\ref{fig:observer})\\
     $+$ PI on odometry \quad $+$ drift observer $\dhat$};
  % ---- bottom row: plant path, right to left ----------------------------------------
  \node[blk, text width=1.5cm] (ekf) at ($(smooth)+(0,-2.6)$) {EKF\\odometry};
  \node[blk, text width=4.2cm, right=0.45cm of ekf] (hull)
    {hull, 3-DOF Fossen\\[1pt] {\scriptsize $M\dot\nuv + C(\nuv_r)\nuv_r + D(\nuv_r)\nuv_r = \bm\tau$}};
  \node[blk, text width=3.0cm, right=0.45cm of hull] (esc)
    {ESC / T200\\[1pt] {\scriptsize deadband 25/400, thrust $\propto e|e|$,\\ spool, dead time}};
  \node[blk, text width=1.8cm, right=0.45cm of esc] (mixer) {force-space\\mixer};
  % ---- signals ------------------------------------------------------------------------
  \draw[sig] (mppi) -- node[above, sub] {$v^*$} (smooth);
  \draw[sig] (smooth) -- (mux);
  \draw[sig] (mux) -- (law);
  \draw[sig] (law.east) -| node[pos=0.75, right, sub] {$e$} (mixer.north);
  \draw[sig] (mixer) -- (esc);
  \draw[sig] (esc) -- node[above, sub] {$\bm\tau$} (hull);
  \draw[sig] (hull) -- node[above, sub] {$\nuv$} (ekf);
  \coordinate (fb) at ($(ekf.north)+(0,0.6)$);
  \draw[thick] (ekf.north) -- (fb);
  \draw[sig] (fb) -| (mppi.south);
  \draw[sig] (fb) -| node[pos=0.3, below, sub] {odometry $\hat\nuv$, $\psi$} ($(law.south)+(-2.0,0)$);
  \fill (fb) circle (1.3pt);
  % ---- what MPPI does not model --------------------------------------------------------
  \node[draw=sA, dashed, rounded corners=4pt, inner sep=5pt,
    fit=(smooth) (mux) (law) (mixer) (esc) (hull) (ekf)] (gap) {};
  \node[anchor=south east, sA, font=\scriptsize] at (gap.north east) {not in MPPI's model};
  % ---- what enters from outside, and where the lag comes from ---------------------------
  \draw[dist] ($(hull.south |- gap.south)+(0,-0.75)$) node[below, note, sB] {waves, current, wind}
    -- (hull.south);
  \draw[dist] ($(ekf.south |- gap.south)+(0,-0.75)$) node[below, note, sB] {noise $\sigma_f$}
    -- (ekf.south);
  \draw[decorate, decoration={brace, amplitude=4pt, mirror}, thick, sD!60!black]
    ($(esc.west |- gap.south)+(0,-0.1)$) -- ($(esc.east |- gap.south)+(0,-0.1)$)
    node[midway, below=4pt, note] {spool + delay: part of $L$};
  \node[note, anchor=north east] at ($(ekf.west |- gap.south)+(-0.1,-0.1)$) {odometry latency:\\ part of $L$};
  \node[note, anchor=north] at ($(gap.south)+(0,-1.25)$)
    {what MPPI meets: the closed-loop velocity lag $\tau$ (the law around the plant)\\
     and the dead time $L$ = spool + delay + odometry latency};
\end{tikzpicture}
\caption{The velocity loop between Nav2 and the water. MPPI plans as if the boat moved
at the commanded velocity; everything inside the dashed box is outside its model. The
velocity law closes a loop around the thrusters and the hull: its feed-forward inverts
the per-direction speed curves at the wanted velocity less the mechanical bias $b$ and
the drift estimate $\dhat$ rotated into the body, and PI feedback on odometry covers the
rest. What MPPI sees of that box is a lag $\tau$, a dead time $L$, the odometry noise
$\sigma_f$ and the disturbances acting on the hull. A hold in \code{command\_mux} sends
neutral thrust, not station keeping.}
\label{fig:loop}
\end{figure}


\subsection{The law}
Let $\bm v^\star = \clip(\bm v^*, -\bm v_\text{max}, \bm v_\text{max})$ be the clipped request,
$\bm b$ the body-fixed mechanical bias, $\dhat$ the world drift estimate and $s_j^{-1}$ the
inverse of axis $j$'s speed curve, $s_j^{-1}(v) = \sgn(w)\sqrt{|w|}$ with $w = d_1 v + d_2 v|v|$.
Per axis,
\[
  e_j = \clip\!\Big(\underbrace{s_j^{-1}\big(v^\star_j - b_j - [R(\psi)^\top\dhat]_j\big)}_{\text{feed-forward}}
        + k_{p,j}\,\epsilon_j + k_{i,j}\,z_j,\; -1,\; 1\Big),
  \qquad \epsilon_j = v^\star_j - v_j ,
\]
with the integral $z_j$ clipped to $\pm z_\text{max}$ (\code{integral\_limit}) and frozen on an
axis pinned at $\pm1$ in the direction of its error (conditional integration). The
defaults are $k_p = (0.8, 0.8, 0.6)$, $k_i = (0.6, 0.6, 0.4)$, $z_\text{max} = 0.3$.

\paragraph{Why invert the curve, and not a gain.} The feed-forward is exact on the model, so
the PI loop only meets what the model gets wrong. A straight-line gain would leave the loop a
speed-dependent error, and a correction through the \emph{secant} gain of a power law
over-corrects: for $y = c\,x^n$ the secant inverse gives $\Delta x = \Delta y\,x/y$, $n$ times
the local inverse $\Delta y/(n c x^{n-1})$. So no measured lag or offset is ``undone'' by
dividing by a secant gain; the curve is inverted, and the residual is left to feedback.

\paragraph{The integral ceiling.} $z_\text{max}$ bounds what the integral can hold. With the
observer off (a stale estimate) it binds only from about 0.45\,m/s of current: goal overshoot
0.39\,m at $z_\text{max} = 0.3$ against 0.27\,m at 0.5 in a 0.45\,m/s current, and no gain at
0.3\,m/s (0.19 against 0.21\,m) \ev{stale\_v1:stale\_il0.3\_c0.45, stale\_il0.5\_c0.45, stale\_il0.3, stale\_il0.5}.
Raising it is optional and only for currents over 0.45\,m/s.
\file{mhseals\_hardware/velocity\_law.py (VelocityLaw.step)}

\subsection{The drift observer}
\label{sec:observer}
Current and wind change from day to day, so the drift is estimated on the boat, in the world
frame where it is constant, and rotated into the body when used. Each odometry sample gives
the part of the velocity the model cannot explain,
\[
  \bm\rho = R(\psi)\,\big(\bm v - \bm b - \bm s(\bm e)\big)_{xy},
\]
and a first-order filter with $\alpha = \Delta t/(\tau_d + \Delta t)$, $\tau_d = 5$\,s, averages it:
$\dhat \leftarrow \dhat + \alpha(\bm\rho - \dhat)$. Three rules decide when it may learn
(Figure~\ref{fig:observer}).

\begin{description}[leftmargin=1em, itemsep=2pt]
\item[Only at rest.] It learns only while every $|v^\star_j| \le 0.1$ (\code{drift\_command\_limit})
  and no effort is saturated. Under way, $\bm\rho$ is mostly the model's own error, which would
  be taken for current and turned with the boat; the estimate is held and the PI covers the
  rest. A held estimate matters: with the observer off, skid in a carved turn in a 0.3\,m/s
  cross-current doubles, 0.054$\to$0.107\,m/s \ev{stale\_v1:observer\_on\_il0.3, stale\_il0.3}.
\item[The sign gate.] Against a dock or a buoy the boat stops while the effort pushes on. The
  residual then reads as a current pushing it \emph{off} the contact, and learning it would
  hold the push and add to it every cycle. A real current that is not yet compensated does the
  opposite: it carries the boat \emph{along} the residual before the loop catches it. So the
  observer integrates the world displacement since the boat last got under way,
  $\bm q \leftarrow \bm q + \Delta t\,R(\psi)\bm v_{xy}$ (on learning samples; reset to 0 whenever
  under way), and lets $|\dhat|$ grow only once
  \[
    \bm q\cdot\frac{\bm\rho}{|\bm\rho|} > q_\text{min} = 0.02\text{\,m}\quad(\code{drift\_carry});
  \]
  otherwise the update may shrink or turn the estimate but $|\dhat|$ is capped at its previous
  size. The gate reads displacement, not speed, because the PI loop soon stops a drifting boat:
  the carry it already had stays on record. The correction is sign-dependent by construction:
  being carried along the residual opens the gate; reaching the dock moving the other way, or
  being pinned, keeps it shut.
\item[The cap.] $|\dhat| \le d_\text{max} = 0.5$\,m/s (\code{drift\_limit}) always: no current
  the boat operates in is faster, and the cap bounds the damage of any path the gate misses. It
  must stay well below the slowest axis's full-effort speed (usually sway), or the feed-forward
  alone could still drive that axis to full effort against a contact.
\end{description}

\noindent The failure these rules remove was seen in crane\_ml's full mission: in the dock
back-out the command stayed under 0.1\,m/s, the observer learned the dock contact as current,
and the effort wound up to $(-0.8, +1.0)$ and pinned the boat to the tine until Nav2 reported
no progress \ev{ma\_v3 seed 2}. The estimate exceeded 0.35\,m/s for 10\,s or more in a
0.15\,m/s current \ev{ma\_v2 seed 0}.

\begin{figure}[tbp]
\centering
\begin{tikzpicture}[>=Stealth, font=\footnotesize,
  blk/.style={draw, thick, rounded corners=2pt, fill=white, align=center, inner sep=3pt,
    minimum height=1.0cm},
  gate/.style={draw, thick, diamond, aspect=2.0, fill=sB!10, align=center, inner sep=0.5pt,
    font=\scriptsize},
  frz/.style={draw=sB, rounded corners=2pt, fill=sB!6, align=center, font=\scriptsize,
    inner sep=3pt},
  sig/.style={->, thick},
  lab/.style={font=\scriptsize}]
  % ---- row 1: when the observer may learn --------------------------------------------
  \node[blk, fill=sA!10] (in) at (0,0) {odometry $v$, $\psi$\\ last effort $e$\\ target $v^*$};
  \node[gate] (way) at (3.55,0) {under way?\\ any $|v^*_i| > 0.1$};
  \node[gate] (sat) at (7.4,0) {saturated?\\ any $|e_i| \ge 0.999$};
  \node[blk] (res) at (12.2,0)
    {residual, world frame\\[1pt]
     $\bm r = R(\psi)\,[\,v - b - \mathrm{speed}(e)\,]_{xy}$\\[1pt]
     carried $\bm q \mathrel{+}= \Delta t\, R(\psi)\,v_{xy}$};
  \draw[sig] (in) -- (way);
  \draw[sig] (way) -- node[above, lab] {no} (sat);
  \draw[sig] (sat) -- node[above, lab] {no} (res);
  \node[frz] (freeze) at (5.475,-1.65) {freeze $\dhat$ \quad (under way also resets $\bm q \leftarrow 0$)};
  \draw[->, sB, thick] (way.south) -- node[left, lab] {yes} (way.south |- freeze.north);
  \draw[->, sB, thick] (sat.south) -- node[right, lab] {yes} (sat.south |- freeze.north);
  % ---- row 2: low-pass, then the sign gate, then the cap ----------------------------
  \node[blk] (lp) at (1.3,-3.9)
    {low-pass\\[1pt] $\dhat \leftarrow \dhat + \alpha(\bm r - \dhat)$\\[1pt]
     $\alpha = \Delta t / (\tau_d + \Delta t)$, \ $\tau_d = 5$\,s};
  \node[gate, aspect=1.7] (sign) at (5.75,-3.9) {carried along $\hat{\bm r}$?\\
    $\bm q \cdot \bm r / |\bm r| > 0.02$\,m};
  \node[blk, fill=sC!10] (grow) at (10.0,-3.25) {$|\dhat| \le 0.5$\,m/s \ may grow};
  \node[blk, fill=sB!6] (keep) at (10.0,-4.55) {$|\dhat| \le \min(|\dhat_{\text{before}}|, 0.5)$\\
    may only shrink or turn};
  \node[blk, fill=sA!10] (out) at (14.0,-3.9) {$R(\psi)^\top \dhat$\\ to the\\ feed-forward};
  \draw[sig] (res.south) -- ++(0,-1.85) -| (lp.north);
  \draw[sig] (lp) -- (sign);
  \draw[sig] (sign.east) -- ++(0.35,0) |- node[pos=0.75, above, lab] {yes} (grow.west);
  \draw[sig] (sign.east) -- ++(0.35,0) |- node[pos=0.75, below, lab] {no} (keep.west);
  \draw[sig] (grow.east) -- ++(0.3,0) |- (out.west);
  \draw[sig] (keep.east) -- ++(0.3,0) |- (out.west);
  % ---- row 3: the side path that bypasses the law -------------------------------------
  \node[frz, font=\footnotesize] (hold) at (1.3,-6.5) {\code{command\_mux} hold or idle\\ or command timeout (0.5\,s)};
  \node[blk] (reset) at (5.3,-6.5) {law reset\\ {\scriptsize integral, effort, target}};
  \node[blk] (zero) at (8.6,-6.5) {zero effort\\ {\scriptsize neutral thrusters}};
  \draw[sig] (hold) -- (reset);
  \draw[sig] (reset) -- (zero);
  \node[lab, align=left, anchor=north west] at ($(hold.south west)+(0,-0.12)$)
    {hold means neutral thrusters, not station keeping; $\dhat$ is kept across the reset};
  % ---- inset: the boat pinned on a dock -------------------------------------------------
  \begin{scope}[shift={(12.9,-6.3)}]
    \fill[pattern=north east lines, pattern color=sD] (1.05,-0.7) rectangle (1.3,0.7);
    \draw[thick] (1.05,-0.7) -- (1.05,0.7);
    \node[lab, rotate=90] at (1.5,0) {dock};
    \draw[thick, fill=sD!15] (-0.75,-0.3) -- (0.75,-0.3) -- (1.05,0) -- (0.75,0.3) -- (-0.75,0.3) -- cycle;
    \draw[->, very thick, sB] (-0.3,0.5) -- node[above, lab] {effort $e$} ++(0.9,0);
    \draw[->, very thick, sA] (0.4,-0.5) -- node[below, lab] {residual $\bm r$} ++(-0.9,0);
    \node[lab, align=center, anchor=north] at (0.1,-1.0) {it arrived moving into the dock,\\ so $\bm q \cdot \hat{\bm r} \le 0$: no growth};
  \end{scope}
\end{tikzpicture}
\caption{The drift observer of the velocity law (\code{mhseals\_hardware} \code{velocity\_law.py}).
It learns only while the boat is asked to hold still and no axis is saturated: the
residual $\bm r$, the body velocity the speed curves do not explain, is turned into the
world frame and low-passed into $\dhat$ with $\tau_d = 5$\,s. The estimate may grow only
once the boat has been \emph{carried} at least 2\,cm along $\bm r$ since it last got under
way; otherwise it may only shrink or turn, and $|\dhat| \le 0.5$\,m/s always. Inset: a
boat pushed against a dock reads its stalled effort as a current pushing it off the
contact, but it reached the contact moving the other way, so the estimate does not grow and the push
does not wind up. Constants: \code{drift\_command\_limit} 0.1, \code{drift\_time\_constant}
5\,s, \code{drift\_carry} 0.02\,m, \code{drift\_limit} 0.5\,m/s.}
\label{fig:observer}
\end{figure}


\paragraph{Hold means neutral.} The controller follows the command mux's latched state: while
its input is anything but auto, manual or RC (a hold, including the operator's \texttt{E}, or idle),
or the state cannot be read, the law is reset and the effort is zero. The command timeout does
the same. The mux sends zeros while it holds; without this the law would station-keep on them,
with the drift feed-forward and the kept integral still pushing. Station
keeping happens only when autonomy explicitly commands zero velocity; a hold never drives the
thrusters to hold position. The remote e-stop is separate: it cuts thruster power in hardware
and is latched until the arm flow re-enables it.

\paragraph{Drift as a gate.} Because the estimate is held under way, a large estimate at the
start of a task is a condition, not a parameter: above 0.35\,m/s precise tasks (docking) do not
start (Section~\ref{sec:field}).
\file{mhseals\_hardware/velocity\_law.py (VelocityLaw.observe), velocity\_controller.py (hold, reset)}
